\section{A Historical Trace of Capacity Utilization Measurement}
\label{sec:literature_review}

This section traces the evolution of capacity utilization measurement from early industrial surveys to contemporary output-gap governance. The metric originally emerged to quantify idle capital during the Great Depression, shifted toward inflation surveillance during the stagflation crisis of the 1970s, and now informs technocratic output-gap benchmarks. Throughout these phases, the unobserved productive ceiling has remained conceptually underidentified. I trace this trajectory through four developments: the standardization of survey measures during the Fordist era, the transition to output-gap governance, the theoretical debates between Sraffian and Kaleckian economists over long-run utilization, and Shaikh's accounting-based framework.

\subsection{From Survey Measures to an Inflation Sentinel: Capacity Utilization during the Fordist Era}
\label{subsec:fordist_survey_statecraft}

Before capacity utilization could serve as a policy benchmark, researchers had to confront a conceptual divergence in its definition. The mass idleness of the Great Depression forced statistical agencies to quantify unused resources, exposing an enduring tension between engineering and economic definitions of productive capacity \citep{Burns1954, KleinDavid1958}. Engineering capacity represents the physical ceiling of machine operation under uninterrupted conditions. Economic capacity, by contrast, defines the profit-maximizing level of output given prevailing wages, input costs, and operating schedules. Early measurement initiatives focused on determining normal operating rates for enterprises amidst the emergence of post-Depression macroeconomic policy.

During the post-war expansion, Alvin Hansen's stagnation thesis warned that chronic underutilization of capital posed an enduring drag on long-run growth \citep{Hansen1936, Hansen1941, Barber1987}. As the neoclassical synthesis codified these concerns into IS-LM demand management \citep{Hicks1937, Samuelson1947, Samuelson1948}, quantitative indicators were required to detect bottlenecks and effective-demand shortfalls. Capacity utilization emerged as a primary metric to render capital idleness legible for countercyclical stabilization.

The McGraw-Hill plant-and-equipment surveys standardized this reporting by asking plant managers to state their operating rates relative to preferred benchmarks. Emerging from the private sector, these surveys served corporate planning and business forecasting before public initiatives took root \citep{BaranSweezy1988, Munroe2007, Butler1958, Phillips1963}. By aggregating firm-level responses into a standardized statistical index, the surveys offered a consistent language for business forecasting. When the Federal Reserve Board subsequently integrated these survey responses into its official Industrial Production index, managerial assessments became an institutional benchmark for macroeconomic surveillance \citep{Board2017}. Alternative physical indicators---such as the workweek of capital \citep{Foss1963, TaubmanGottschalk1971}---never secured comparable bureaucratic adoption.

The stagflation crisis of the 1970s dismantled the post-war demand-stabilization framework. Faced with falling profit rates and rising inflation, policymakers repurposed the capacity index: it functioned less as an indicator of effective-demand shortfalls to be expanded, and more as a sentinel for inflation control \citep{ShaikhManiatisPetralias1999}. High capacity utilization was treated as a warning of impending price pressures, prompting monetary tightening. Academic and policy discussions increasingly evaluated capacity measures by their ability to predict inflation \citep{KleinEtAl1973}. In these debates, Alan Greenspan argued that capacity measures must reflect systemic shortages rather than administrative averages, asserting that ``if the numbers do not indicate that the economy is currently at capacity, then I suggest that the numbers are wrong, and that the correct concept of capacity must reflect the situation'' \citep[p.~758]{KleinEtAl1973}. This shift established the conceptual foundation for modern output-gap governance.

\subsection{From Capacity Utilization to Output-Gap Governance}
\label{subsec:output_gap_governance}

During the late twentieth century, the growth of international trade and financial integration reduced the policy primacy of national manufacturing surveys. Central banks and international agencies turned toward portable, rule-based indicators to monitor macroeconomic conditions \citep{Jessop2002, Jessop2003, Konings2007}. This shift elevated the ``output gap''---the difference between actual and potential output---as the core operational anchor for monetary frameworks, including inflation targeting and Taylor rules \citep{Taylor1993, Johnson2024}. Under the New Consensus Macroeconomics, potential output was defined as the non-accelerating inflation limit of economic activity, often linked to NAIRU benchmarks \citep{SaadFilho2018}. Yet potential output remained an unobservable latent variable whose trajectory depends heavily on statistical filtering choices \citep{HerndonAshPollin2014, AshBasuDube2017}.

Because potential output cannot be observed directly, its estimation depends on specific modeling assumptions. Mainstream practice typically relies on statistical filters that smooth observed output \citep{HodrickPrescott1997, BaxterKing1999, BeveridgeNelson1981, Hamilton2018}, aggregate production functions \citep{CBO2004PotentialGDP, Alichi2015}, or structural vector autoregressions \citep{BlanchardQuah1989}. These methodologies treat potential output as a frictionless trend or an engineering trajectory. In political economy, by contrast, the conversion of productive equipment into realized output is understood as institutionally mediated. At the plant level, work schedules, line speeds, and labor discipline govern how intensively physical equipment can run. In the macroeconomy, the realization of potential output interacts with distributive conflict, reserve-army dynamics, and profitability. When low unemployment strengthens labor's bargaining position, resulting profit-squeeze pressures may lead firms to idle capacity rather than expand production. Distinguishing between physical equipment and institutional conditions is therefore central to measuring the capacity ceiling.

The 2008 Global Financial Crisis exposed the practical limits of this filtering approach. Post-crisis revisions routinely altered the sign of estimated output gaps, and downward revisions to potential output were used to justify fiscal austerity in the European periphery \citep{Heimberger2017, Heimberger2020}. Although international institutions documented substantial hysteresis and revision volatility \citep{CerraSaxena2017, CerraFatasSaxena2023}, surveillance frameworks continued to rely on output-gap calculations for real-time policy guidance \citep{AastveitTrovik2014, Foster2024}. This revision-prone setting provides the institutional context for heterodox alternatives to capacity measurement.

\subsection{Capacity-Utilization Debates in Political Economy}
\label{subsec:pe_utilization_debates}

Within political economy, the ``capacity utilization controversy'' divides scholars along a central question: is the long-run normal rate of utilization endogenous to demand, or is it anchored in cost structures and classical competition? The Neo-Kaleckian tradition treats the normal rate as endogenous, arguing that firms adjust long-run operating rates in response to demand shocks, shift choices, or historical hysteresis \citep{Lavoie1995, Nikiforos2013, Setterfield2020}. In contrast, Sraffian economists argue that cost minimization and classical competition anchor the normal rate of utilization in the long run, with capital investment adjusting capacity to meet changes in effective demand \citep{Ciccone1986, Kurz1986}.

Both positions face empirical and conceptual challenges when identifying the productive ceiling. In the Sraffian supermultiplier framework, convergence to a normal rate requires an autonomous component of demand, a condition that can generate tensions with the debt dynamics of stock-flow consistent models \citep{Nikiforos2018}. Conversely, the Neo-Kaleckian hypothesis of endogenous utilization requires capacity utilization to exhibit long-run non-stationarity. Empirical tests of this hypothesis remain contested: \citet{GahnGonzalez2020, GahnGonzalez2022} argue that capacity utilization is stationary and mean-reverting, whereas \citet{Nikiforos2020} shows that unit-root inferences depend on sample selection and deterministic specifications.

Furthermore, by focusing primarily on manufacturing, both the FRB and alternative survey series miss the broader structural shifts of modern, service-dominated economies \citep{Haluska2023}. Evaluating the long-run dynamics of capacity formation therefore motivates an empirical strategy that departs from survey indices and statistical smoothing.

\subsection{Shaikh's Structural Identification of Capacity Utilization}
\label{subsec:shaikh_structural_measurement}

Rather than relying on survey measures or statistical filters that impose symmetric, mean-reverting cycles around a smooth trend, \citet{Shaikh2016} anchors the capacity ceiling in capital accumulation and profitability accounting. Standard filtering techniques treat expansions and contractions as mirror images, obscuring the cyclical asymmetry between prolonged booms and rapid contractions---a critique shared by neo-structuralist analyses of persistent recessive gaps and hysteresis \citep{FfrenchDavis2010, FfrenchDavis2012}. Shaikh's framework addresses this limitation by grounding potential output in the dynamics of the capital stock and the maximum potential profit rate. Using a Generalized Perpetual Inventory Method (GPIM) to construct gross capital stock, he identifies capacity utilization as the deviation of actual output from this estimated ceiling, retaining the historical trace of accumulation and scrapping.

This accounting-based approach offers two distinct methodological advantages: it relies on reproducible national accounts data and distinguishes net capital stock (which tracks valuation) from gross capital stock (which reflects physical assets in operation). However, treating capacity utilization as a long-run residual assumes that output and capital share a stable cointegrating relationship. If this empirical relation is sensitive to specification or structural breaks, the resulting utilization index inherits that instability. In this paper, I critically replicate Shaikh's cointegration framework to assess its empirical stability across lag choices, deterministic terms, and sample periods, establishing the empirical foundation for the distributional extension that follows.
